\begin{solapartat}
\punts{0,2} $v_G = \dfrac{2\pi R_G}{T_G}$

\punts{0,1} $T_G = 0,428\ \text{dies}\times\dfrac{24\ \mathrm{h}}{1\ \text{dia}}\times\dfrac{60\ \mathrm{min}}{1\ \mathrm{h}}\times\dfrac{60\ \mathrm{s}}{1\ \mathrm{min}} = 3,70\times 10^{4}\ \mathrm{s}$

\punts{0,2} $v_G = \dfrac{2\pi\times 6,20\times 10^{7}}{3,70\times 10^{4}} = 1,05\times 10^{4}\ \mathrm{m/s}$

\punts{0,3} La força d'atracció gravitatòria és igual a la força centrípeta necessària perquè el satèl·lit giri en la seva òrbita.
\[ F_{\textit{gravitatòria}} = F_{\textit{centrípeta}} \Rightarrow G\frac{M_N m_G}{{R_G}^2} = \frac{m_G v_G^2}{R_G} \]

\punts{0,2} $M_N = \dfrac{R_G v_G^2}{G} = \dfrac{6,20\times 10^{7}\times\left(1,05\times 10^{4}\right)^2}{6,67\times 10^{-11}} = 1,02\times 10^{26}\ \mathrm{kg}$
\end{solapartat}

\begin{solapartat}
\punts{0,5} $g_N = \dfrac{GM_N}{{R_N}^2}$

\punts{0,5} $g_N = 11,2\ \mathrm{m/s^2}$ ó $(\mathrm{N/kg})$
\end{solapartat}
